% ============================================================
% Phân tích use cases: Hoàn trả + Tranh chấp
% ============================================================

\subsection{Đặc tả use case tạo hoàn trả (UC18)}
\label{uc:create-refund}

\begin{table}[htbp]
\centering
\caption{Đặc tả UC18 --- Tạo yêu cầu hoàn trả}
\label{tab:uc-create-refund}
\small
\begin{tabular}{@{} p{3cm} p{10.5cm} @{}}
\toprule
Mục đích & Buyer tạo yêu cầu hoàn trả cho đơn hàng đã thanh toán thành công \\
Actor & Buyer; liên quan: Seller (thông báo), Analytic (ghi nhận refund\_requested) \\
Tiền ĐK & Đã đăng nhập; order thuộc buyer, payment=Success; chưa có refund Pending/Processing \\
Hậu ĐK & Refund status=Pending; bằng chứng liên kết qua Common resource; thông báo seller \\
Luồng chính & Chọn order $\rightarrow$ nhập lý do, phương thức (PickUp/DropOff), bằng chứng $\rightarrow$ tạo refund $\rightarrow$ thông báo \\
Ngoại lệ & Chưa thanh toán hoặc đã có refund đang xử lý $\rightarrow$ từ chối \\
\bottomrule
\end{tabular}
\end{table}

\subsection{Đặc tả use case duyệt hoàn trả (UC21)}
\label{uc:confirm-refund}

\begin{table}[htbp]
\centering
\caption{Đặc tả UC21 --- Duyệt yêu cầu hoàn trả}
\label{tab:uc-confirm-refund}
\small
\begin{tabular}{@{} p{3cm} p{10.5cm} @{}}
\toprule
Mục đích & Seller chấp nhận hoặc từ chối yêu cầu hoàn trả \\
Actor & Seller; liên quan: Payment Gateway (auto-refund), Transport Provider \\
Tiền ĐK & Refund status=Pending; thuộc đơn hàng của seller \\
Hậu ĐK & Chấp nhận: status=Processing, auto-refund nếu có provider\_charge\_id, tạo transport trả hàng (PickUp). Từ chối: status=Cancelled \\
Luồng chính & Xem lý do + bằng chứng $\rightarrow$ chấp nhận (trigger auto-refund) hoặc từ chối $\rightarrow$ thông báo buyer \\
\bottomrule
\end{tabular}
\end{table}

\subsection{Đặc tả use case tranh chấp (UC22)}
\label{uc:open-dispute}

\begin{table}[htbp]
\centering
\caption{Đặc tả UC22 --- Mở tranh chấp hoàn trả}
\label{tab:uc-open-dispute}
\small
\begin{tabular}{@{} p{3cm} p{10.5cm} @{}}
\toprule
Mục đích & Buyer hoặc Seller mở tranh chấp khi không đồng ý quyết định hoàn trả \\
Actor & Buyer hoặc Seller; liên quan: Admin (giải quyết --- UC23) \\
Tiền ĐK & Refund tồn tại; bên mở tranh chấp liên quan đến đơn hàng \\
Hậu ĐK & Tạo refund\_dispute status=Pending, lưu issued\_by\_id và lý do \\
Luồng chính & Chọn ``Mở tranh chấp'' $\rightarrow$ nhập lý do $\rightarrow$ tạo dispute $\rightarrow$ chờ Admin xử lý \\
\bottomrule
\end{tabular}
\end{table}

\subsection{Trạng thái hoàn trả và tranh chấp}

\begin{table}[htbp]
\centering
\caption{Chuyển trạng thái --- Hoàn trả và tranh chấp}
\label{tab:refund-states}
\small
\begin{tabular}{@{} l l l l @{}}
\toprule
\rowcolor{headerblue}
\thead{Thực thể} & \thead{Từ} & \thead{Đến} & \thead{Sự kiện} \\
\midrule
Refund & --- & Pending & CreateBuyerRefund \\
Refund & Pending & Processing / Cancelled & ConfirmSellerRefund (chấp nhận / từ chối) \\
Refund & Processing & Success / Failed & Kết quả hoàn tiền \\
Dispute & --- & Pending & Buyer/Seller mở tranh chấp \\
Dispute & Pending & Processing $\rightarrow$ Success / Cancelled & Admin xem xét / hai bên tự thỏa thuận \\
\bottomrule
\end{tabular}
\end{table}

\begin{figure}[htbp]
\centering
\resizebox{\textwidth}{!}{%
\begin{sequencediagram}
\newinst{buyer}{Khách hàng}
\newinst[3]{shop}{ShopNexus}
\newinst[3]{seller}{Người bán}
\newinst[3]{admin}{Quản trị viên}

\begin{call}{buyer}{createRefund(orderID, method, reason)}{shop}{refund}
\end{call}

\begin{messcall}{shop}{notify(refundRequested)}{seller}
\end{messcall}

\begin{call}{seller}{confirmRefund(refundID)}{shop}{refund}
\end{call}

\begin{messcall}{shop}{autoRefund(paymentID)}{shop}
\end{messcall}

\begin{messcall}{shop}{notify(refundApproved)}{buyer}
\end{messcall}

\postlevel
\begin{sdblock}{alt}{Seller từ chối}

\begin{call}{buyer}{createDispute(refundID, reason)}{shop}{dispute}
\end{call}

\begin{messcall}{shop}{notify(disputeCreated)}{admin}
\end{messcall}

\begin{call}{admin}{resolveDispute(disputeID, status)}{shop}{result}
\end{call}

\end{sdblock}

\end{sequencediagram}
}%
\caption{Lược đồ tuần tự --- Luồng hoàn trả và tranh chấp}
\label{fig:refund-sequence}
\end{figure}

\subsection{Thành phần nghiệp vụ}

\textbf{Refund}: ListBuyerRefunds, CreateBuyerRefund, UpdateBuyerRefund, CancelBuyerRefund, ConfirmSellerRefund. \textbf{Dispute}: CreateRefundDispute, ResolveRefundDispute. Helper: Payment Gateway (auto-refund), Transport Provider (vận đơn trả hàng PickUp).
